% ---------------------------------------------------------------------
% Immunology paragraph for the Introduction of the CNS submission.
% ~500 words. Contributed by Immunology agent (agent-mozus3vv) for
% task t-mozv8ou4 (follow-up: Biological Science asked for this
% explicitly as workspace/handoffs/immuno_intro_para.tex).
% Target: sections/01_introduction.tex, probably after the general
% motivation paragraph on integration + overcorrection and before
% the paragraph that introduces our framework (LRS / RareShield).
% ---------------------------------------------------------------------

The stakes of this failure mode are not uniformly distributed across
the cell-biology landscape; they are most concentrated in the immune
compartment, which is also the compartment where the majority of
integration-atlas effort now accumulates. Three structural properties
of immunology make unannotated rare subpopulations especially
vulnerable to batch-integration overcorrection. First, the immune
system is a lineage continuum: committed lineages pass through
progenitor, activation, and terminal states whose rare intermediates
(pre-DC, transitional B, progenitor-exhausted CD8) sit at the joints
of an explicitly hierarchical tree \citep{villani2017single,
cancro2020abcs}. Because integration methods penalize between-batch
dispersion in a shared latent space, these joint populations --- the
populations that are rare precisely because they are transitional ---
are the populations with the highest expected displacement under
batch alignment. Second, an immune cell's transcriptome is tissue-,
antigen-, and context-dependent to a degree rarely matched elsewhere:
a circulating memory CD8 T cell and its tissue-resident counterpart
share at most 60\% of their differentially expressed genes, and
mucosal-associated invariant T cells (MAIT) in blood, gut, and lung
cluster as effectively distinct populations despite identical TCR
identity \citep{provine2020mait, dominguez2022crosstissue}. Standard
batch covariates (patient, chemistry, sequencing centre) collapse
precisely those axes, rendering biologically informative structure
indistinguishable from technical noise. Third, prevalence itself is
informative: the clinical readout in immunology is often the fold
change in a rare population's abundance across conditions --- MAIT
collapse in severe COVID-19, age-associated B-cell expansion in
systemic lupus erythematosus, CD103$^+$CD39$^+$ tumour-reactive
tissue-resident memory CD8 ratios between checkpoint-blockade
responders and non-responders. Integration methods that normalize
cell-type prevalence across batches to ``correct'' sampling
imbalance obliterate exactly the signal the downstream analysis is
designed to measure \citep{duhen2018coexpression,
caushi2021transcriptional, vandamme2021depletion}.

The historical record amplifies the concern. Nearly every
mechanistic advance in the past decade of immunology began as a
small, ambiguously-identified cluster that predated its nameable
identity. Plasmacytoid dendritic cells, at $<\!0.5\%$ of peripheral
blood and transcriptomically adjacent to B cells, would not have
established the type-I interferon axis of antiviral immunity had
they been absorbed into their abundant neighbour during integration;
the \textit{AXL}$^+$\textit{SIGLEC6}$^+$ ``AS-DC'' population (DC5),
at $\sim\!0.05\%$ of PBMC, rewrote the classification of the
conventional DC lineage only after Villani \etal\ isolated it as a
distinct cluster from several thousand cells
\citep{villani2017single}; tumour-reactive \textit{CD103$^+$CD39$^+$}
tissue-resident memory T cells at $<\!1\%$ of tumour cells now drive
adoptive-cell-therapy patient selection in multiple trials
\citep{duhen2018coexpression}. In every case, the cluster existed
before the biology was understood, and in every case, an integration
pipeline that silently merged the cluster into its abundant
neighbour would have delayed or extinguished the discovery. To the
degree that pan-cancer and pan-tissue atlases now dominate immune
biology \citep{kang2024integrated, cheng2021pancancer,
zheng2021pancancer}, integration is no longer a pre-processing step;
it is the step at which future rare-population discoveries are
either protected or silently lost. A detection framework that
operates without relying on prior annotation, and an integration
strategy that is rare-aware by construction, are therefore not
refinements of an existing toolkit but the prerequisites for the
next generation of immunological discovery.

% ---------------------------------------------------------------------
% Requires: villani2017single, cancro2020abcs, provine2020mait,
% dominguez2022crosstissue, duhen2018coexpression,
% caushi2021transcriptional, vandamme2021depletion, kang2024integrated,
% cheng2021pancancer, zheng2021pancancer. All in
% workspace/handoffs/bib_additions_immunology.bib.
% ---------------------------------------------------------------------
